% Generated draft: requires mathematical and source review.
\documentclass[11pt]{article}
\usepackage[margin=25mm]{geometry}
\usepackage{fontspec}
\setmainfont{Latin Modern Roman}
\usepackage{amsmath,amssymb,mathtools,unicode-math}
\setmathfont{Latin Modern Math}
\usepackage{enumitem,needspace,xurl,hyperref}
\setlength{\parindent}{0pt}
\setlength{\parskip}{5pt}
\setlength{\emergencystretch}{3em}
\newcommand{\N}{\mathbb{N}}
\newcommand{\Z}{\mathbb{Z}}
\newcommand{\Q}{\mathbb{Q}}
\newcommand{\R}{\mathbb{R}}
\newcommand{\C}{\mathbb{C}}
\newcommand{\F}{\mathbb{F}}
\newcommand{\A}{\mathbb{A}}
\newcommand{\PP}{\mathbb{P}}
\newcommand{\E}{\mathbb{E}}
\newcommand{\eps}{\varepsilon}
\newcommand{\dd}{\,\mathrm{d}}
\begin{document}
\section*{UnsolvedMath ID: 7000023}
\textbf{Source problem code:} catalogue entry
\subsection{Definitions and mathematical statement}
Let $M$ be a compact connected minimal surface in $\mathbb R^3$ whose boundary consists of two convex planar curves. A convex planar curve is a simple closed curve bounding a convex region in its plane; each curve may lie in its own plane. Minimality means vanishing mean-curvature vector in the surface interior, not necessarily area minimization.

Must $M$ be homeomorphic to the annulus $S^1\times[0,1]$? In an immersion formulation, this means a compact connected surface with two boundary components and an immersion $f:M\to\mathbb R^3$ whose two boundary curves have the stated convexity property. In an embedded formulation one additionally requires $f$ to be an embedding. The source does not say which category is intended.

No assumption that the two boundary planes are parallel, and no strict-convexity, stability, or area-minimizing condition, is added. These stronger hypotheses would describe special cases rather than the literal question. The supplied summary refers to Nitsche/Plateau-type results and catenoid examples but provides no precise theorem establishing this general topological conclusion.
\subsection{Short English statement}
Must a compact connected minimal surface with exactly two convex planar boundary curves be an annulus? Preserve the general boundary-plane geometry and clarify whether the surface is required to be embedded.
\subsection{Sources}
\begin{itemize}
\item Mohammad Ghomi, \emph{Open Problems in Geometry of Curves and Surfaces} (2019), Problem 6.1, PDF page 14, \url{https://people.math.gatech.edu/~ghomi/Papers/op.pdf}; accessed 2026-07-29, PDF-text-normalized extraction.
\item The supplied research summary mentions Nitsche/Plateau-type results and catenoid configurations without full bibliographic citations or precise hypotheses. Its special-case descriptions are not used as a proof of the general statement.
\item Local registry \texttt{data/unsolvedmath/problems.json}, permanent UnsolvedMath ID $7000023$, source code catalogue entry, \emph{Geometry of Curves and Surfaces --- Problem 6.1}. Dataset attribution: Ulam AI and the UnsolvedMath Contributors.
\end{itemize}
\end{document}
